\documentclass[12pt]{article}
\usepackage[a4paper,margin=2.5cm]{geometry}
\usepackage[T1]{fontenc}
\usepackage[utf8]{inputenc}
\usepackage{lmodern}
\usepackage{microtype}
\usepackage{amsmath}
\usepackage{xcolor}
\usepackage{listings}
\usepackage[hidelinks]{hyperref}
\definecolor{codegray}{RGB}{245,245,245}
\lstset{basicstyle=\ttfamily\small,backgroundcolor=\color{codegray},frame=single,
  breaklines=true,showstringspaces=false,columns=fullflexible,keepspaces=true}
\setlength{\parindent}{0pt}
\setlength{\parskip}{0.55em}

\begin{document}
\pagenumbering{gobble}
\begin{titlepage}
\centering
{\large Ministerul Educa\c{t}iei \c{s}i Cercet\u{a}rii al Republicii Moldova\par}
\vspace{0.45cm}
{\large Universitatea Tehnic\u{a} a Moldovei\par}
\vspace{0.45cm}
{\large Facultatea Calculatoare, Informatic\u{a} \c{s}i Microelectronic\u{a}\par}
\vspace{0.45cm}
{\large Departamentul Ingineria Software \c{s}i Automatic\u{a}\par}
\vspace{3cm}
{\Large\bfseries Report on Laboratory Work No. 6\par}
\vspace{0.55cm}
{\LARGE\bfseries Inside the xv6 File System\par}
\vspace{0.35cm}
{\large Operating Systems\par}
\vfill
\begin{flushleft}
\textbf{Done by:} st. gr. FAF-241, Ca\c{t} Chiril\\[0.8cm]
\textbf{Checked by:} asist. univ. Rei\c{t}man Patricia
\end{flushleft}
\vfill
{Chi\c{s}in\u{a}u -- 2026\par}
\end{titlepage}

\pagenumbering{roman}
\tableofcontents
\clearpage
\pagenumbering{arabic}

\section{Purpose and environment}
I inspected the xv6 file system's directory entries, inodes, hard links, and disk
layout, then ran two user programs: \texttt{mystat}, which reports metadata for
paths, and \texttt{find}, which walks directories recursively. I used the
\texttt{os/xv6-riscv} tree with the local RISC-V toolchain and QEMU on
26 October 2026. I rebuilt the file system image with
\texttt{make clean \&\& make fs.img}. The selected command output below is
also recorded in \texttt{report\_files/notes.md}. Inode numbers depend on the
programs present in the image.

\section{Directory entries and file types}
\texttt{kernel/stat.h} defines type 1 as a directory (\texttt{T\_DIR}), type 2
as a regular file (\texttt{T\_FILE}), and type 3 as a device
(\texttt{T\_DEVICE}). The root has inode 1, as defined by \texttt{ROOTINO}.
Its \texttt{.} and \texttt{..} both refer to inode 1 because the root has no
parent directory above it.

\begin{lstlisting}
$ ls .
.              1 1 1024
..             1 1 1024
README         2 2 2441
... (other built-in programs)
console        3 27 0
$ mkdir a
$ ls a
.              1 28 32
..             1 1 1024
\end{lstlisting}

An empty \texttt{a} has two entries, \texttt{.} (inode 28) and \texttt{..}
(inode 1). Its 32-byte size follows from two 16-byte \texttt{dirent}
records. Each record stores a 2-byte inode number and a 14-byte name.

\clearpage
\section{Metadata and hard links}
\texttt{mystat} calls \texttt{stat()} for each argument. It prints the file
type, inode number, link count, and size; an invalid path is reported on
standard error without stopping subsequent arguments.

\begin{lstlisting}
$ mystat README . console
README: file, inode 2, links 1, size 2441 bytes
.: directory, inode 1, links 2, size 1024 bytes
console: device, inode 27, links 1, size 0 bytes
$ ln README r2
$ mystat README r2
README: file, inode 2, links 2, size 2441 bytes
r2: file, inode 2, links 2, size 2441 bytes
$ rm README
$ mystat r2
r2: file, inode 2, links 1, size 2441 bytes
$ cat r2 | wc
48 336 2441
\end{lstlisting}

Creating \texttt{r2} added a second directory entry for inode 2, so the
link count became 2 while inode and size stayed the same. Removing
\texttt{README} reduced the count to 1. The content remained readable through
\texttt{r2}. The root's link count was 2 here because \texttt{mkdir a} had
already added a child directory.

In \texttt{kernel/sysfile.c}, \texttt{sys\_unlink} performs
\lstinline!ip->nlink--;! and \texttt{iupdate(ip);}. When the count becomes zero
and the final in-memory reference is released, \texttt{iput()} in
\texttt{kernel/fs.c} invokes \texttt{itrunc(ip)} to free the data blocks.
The on-disk \texttt{struct dinode} contains \texttt{type}, \texttt{major},
\texttt{minor}, \texttt{nlink}, \texttt{size}, and
\texttt{addrs[NDIRECT + 1]}: twelve direct block addresses and one indirect
block address. It has no permissions or symbolic-link field.

\clearpage
\section{File system limits}
The constants in \texttt{kernel/fs.h} are \texttt{BSIZE=1024} and
\texttt{NDIRECT=12}. The indirect block holds
\texttt{NINDIRECT=BSIZE/sizeof(uint)=256} addresses. The directory name
limit is \texttt{DIRSIZ=14}, so a longer name is truncated to 14 bytes:

\begin{lstlisting}
$ mkdir averyveryverylongname
$ ls .
averyveryveryl 1 29 32
\end{lstlisting}

The largest file is
\[
  \left(12+\frac{1024}{4}\right)\times1024
  = 274{,}432\text{ bytes} = 268\text{ KiB}.
\]
This uses twelve direct blocks and 256 block addresses in one indirect block.
The image builder printed:

\begin{lstlisting}
nmeta 47 (boot, super, log blocks 31, inode blocks 13, bitmap blocks 1) blocks 1953 total 2000
\end{lstlisting}

The disk has 2000 blocks in total: 31 for the log, 13 for inodes, 1953 for
data, and one each for the boot block, superblock, and bitmap. Thus the
metadata count is $2+31+13+1=47$.

\clearpage
\section{Recursive find}
\texttt{find} checks that its first argument is a directory, reads each
\texttt{dirent}, ignores free entries and \texttt{.}/\texttt{..}, compares the
last name component with \texttt{strcmp}, and recursively visits child
directories. It closes each opened directory. Its path buffer is checked
before appending a child name.

\begin{lstlisting}
$ mkdir a/b
$ echo x > a/b/r2
$ echo y > a/r2
$ find . r2
./a/b/r2
./a/r2
./r2
$ find a b
a/b
\end{lstlisting}

The three paths match by last component. The root-level \texttt{r2} is the
remaining hard link to the original \texttt{README}; the two files created
with \texttt{echo} are separate files. An additional link check confirmed
that \texttt{find} reports a linked name and \texttt{mystat} identifies its
inode:

\begin{lstlisting}
$ ln cat r3
$ find . r3
./r3
$ mystat cat r3
cat: file, inode 7, links 2, size 37344 bytes
r3: file, inode 7, links 2, size 37344 bytes
\end{lstlisting}

\section{Conclusion}
The experiments showed that directory names map to inodes, hard links share
an inode and its data, and removal frees file data only after no links and no
open references remain. The image layout and file-size limit follow directly
from xv6's constants. Both new programs built and ran in QEMU; \texttt{find}
found nested matches and \texttt{mystat} reflected the link-count changes.

\end{document}
